% Generated draft: requires mathematical and source review.
\documentclass[11pt]{article}
\usepackage[margin=25mm]{geometry}
\usepackage{fontspec}
\setmainfont{Latin Modern Roman}
\usepackage{amsmath,amssymb,mathtools,unicode-math}
\setmathfont{Latin Modern Math}
\usepackage{enumitem,needspace,xurl,hyperref}
\setlength{\parindent}{0pt}
\setlength{\parskip}{5pt}
\setlength{\emergencystretch}{3em}
\newcommand{\N}{\mathbb{N}}
\newcommand{\Z}{\mathbb{Z}}
\newcommand{\Q}{\mathbb{Q}}
\newcommand{\R}{\mathbb{R}}
\newcommand{\C}{\mathbb{C}}
\newcommand{\F}{\mathbb{F}}
\newcommand{\A}{\mathbb{A}}
\newcommand{\PP}{\mathbb{P}}
\newcommand{\E}{\mathbb{E}}
\newcommand{\eps}{\varepsilon}
\newcommand{\dd}{\,\mathrm{d}}
\begin{document}
\section*{UnsolvedMath ID: 701036}
\textbf{Source problem code:} catalogue entry
\subsection{Definitions and mathematical statement}
Work in the language of groups, with multiplication, inversion, and identity element \(1\). A group word \(w(x_1,\ldots,x_n)\) is an identity of a group \(G\) if \(w(g_1,\ldots,g_n)=1\) for every \((g_1,\ldots,g_n)\in G^n\). Write \(\operatorname{Id}(G)\) for the set of all such identities, over all finite numbers of variables. Every equation between group words can equivalently be written in this form.

A finite basis for the identities of \(G\) is a finite set \(B_G\subseteq\operatorname{Id}(G)\) such that, for every group \(H\),
\[
 H\models B_G
 \quad\Longleftrightarrow\quad
 H\models\operatorname{Id}(G).
\]
Here \(H\models S\) means that every identity in \(S\) holds under every assignment of its variables in \(H\). Thus all identities of \(G\) must follow from \(B_G\) together with the group axioms.

A group \(G\) is abelian-by-nilpotent, abelian-by-finite, or abelian-by-(finite nilpotent), respectively, if it has an abelian normal subgroup \(A\trianglelefteq G\) for which \(G/A\) is, respectively, nilpotent, finite, or both finite and nilpotent. For a group \(Q\), define \(\gamma_1(Q)=Q\) and \(\gamma_{j+1}(Q)=[\gamma_j(Q),Q]\), where \([U,V]\) is generated by the commutators \(u^{-1}v^{-1}uv\) with \(u\in U\) and \(v\in V\). The group \(Q\) is nilpotent if \(\gamma_{c+1}(Q)=\{1\}\) for some integer \(c\geq 0\).

The source presents the following as three ``well-known problems.'' Determine, separately, whether each assertion is true:
\begin{enumerate}
\item[(a)] Every abelian-by-nilpotent group \(G\) has a finite basis for \(\operatorname{Id}(G)\).
\item[(b)] Every abelian-by-finite group \(G\) has a finite basis for \(\operatorname{Id}(G)\).
\item[(c)] Every abelian-by-(finite nilpotent) group \(G\) has a finite basis for \(\operatorname{Id}(G)\).
\end{enumerate}
The groups \(G\) and the abelian normal subgroups \(A\) are not assumed finite or finitely generated. The finite basis may depend on \(G\); neither a common basis nor a uniform bound on its size is requested. The quotient order and nilpotency class, where applicable, may also depend on \(G\). The supplied registry labels the problem open; this label is not an independently verified assessment.
\subsection{Short English statement}
Does every group with an abelian normal subgroup and a nilpotent quotient have finitely based group identities? Ask the same question separately when the quotient is finite, and when it is both finite and nilpotent. Each individual group may have its own finite basis.
\subsection{Sources}
\begin{itemize}
\item Local UnsolvedMath registry, permanent ID \(701036\), exact source problem code \texttt{catalogue entry}, title ``Kourovka Notebook Problem 14.51.'' The supplied record labels its status \texttt{open}.
\item The supplied attribution is to A. N. Krasil'nikov, proposed in 1999. The background identifies the source as \emph{The Kourovka Notebook (Unsolved Problems in Group Theory)}, 14th issue (1999), Problem 14.51, as included in the 21st issue edition (2026), arXiv:1401.0300v48: \url{https://arxiv.org/abs/1401.0300v48}. The supplied collection description names E. I. Khukhro and V. D. Mazurov as editors.
\item The source statement calls these ``well-known problems.'' The cited edition and the supplied open-status label have not been independently checked for this draft.
\end{itemize}
\end{document}
